\documentclass[11pt]{article}
\usepackage[margin=1in]{geometry}
\usepackage{amsmath,amssymb,amsthm}
\newtheorem{theorem}{Theorem}
\newtheorem{lemma}{Lemma}
\newtheorem{remark}{Remark}
\title{Exposed-Wall Geometry Immediately Above One Radian}
\author{}\date{October 5, 2026 --- paper proof}
\begin{document}\maketitle
Let $w=1+\delta$ and let $K_\delta$ be any global maximizing cap. Use the active supports and differences $p_\delta=p_\delta^*+F_\delta$, $k_\delta=k_\delta^*+G_\delta$ from \texttt{transition-asymptotics.tex}. Put
\[
 \alpha_\delta=p_\delta(0)-1,\qquad
 W_\delta(t)=\frac{p_\delta(t)-1}{\cos t},\qquad
 Z_\delta(s)=\frac{W_\delta(\delta s)-\alpha_\delta}{\delta^3}.
\]
The right tail is the niche portion with $x>\alpha_\delta$. Reflection gives the left-tail analogues without assuming the maximizing cap is symmetric.

\begin{lemma}[Only angles of order $\delta$ can produce the right tail]
There is a constant $S<\infty$, independent of the maximizing cap for small $\delta$, such that $W_\delta(t)>\alpha_\delta$ implies $t<S\delta$. Moreover
\[
 Z_\delta\longrightarrow V_*(s):=q_*(s)+s^2/2-s^3/6
\]
uniformly on every fixed compact $s$ interval, including zero.
\end{lemma}
\begin{proof}
The support-profile theorem gives uniform convergence on compact intervals of $\delta^{-3}(F_\delta(\delta s)-F_\delta(0))$ to $q_*$. The exact intercept identity and $\|F_\delta\|_\infty=O(\delta^{5/2})$ then give the stated uniform convergence, also at zero.

To exclude escape of the maximizing angle, Cauchy--Schwarz gives
$|F_\delta(t)-F_\delta(0)|\le C\delta^{5/2}\sqrt t$. For $t$ in a fixed sufficiently small neighborhood of zero, the relaxed intercept identity implies
\[
 W_\delta(t)-\alpha_\delta
 \le C\delta t^2-c t^3+C\delta^{5/2}\sqrt t+C\delta^{5/2}t^2
\]
for constants $C,c>0$. Substituting $t=\delta s$ and dividing by $\delta^3$ bounds this by $Cs^2-cs^3+C\sqrt s+C\delta^{3/2}s^2$, negative for all sufficiently large fixed $s$. Away from zero the angle-one relaxed intercept difference is $(\sin t-t)/\cos t<0$, uniformly on closed subintervals of the varying domain. Uniform support convergence keeps it negative there. Together these bounds give a single constant $S$.
\end{proof}

\begin{theorem}[Floor extent, exposed angles, and the limiting curve]
The rightmost floor intercept of the niche satisfies
\[
 \sup_{0<t<w}W_\delta(t)-\alpha_\delta=\frac13\delta^3+o(\delta^3).
\]
Every angle maximizing this intercept satisfies $t/\delta\to2$.

For $X\in(0,1/3)$ let $s(X)\in(3/2,2)$ be the unique solution of
\[
 X=\frac{s^2(2s-3)}{12}.
\]
At $x=\alpha_\delta+\delta^3X$, the niche height, divided by $\delta^2$, converges to $s(X)/2-s(X)^2/4$. The convergence is uniform on compact subintervals of $(0,1/3)$. Any angle realizing the upper envelope at those abscissas satisfies $t/\delta\to s(X)$, uniformly in this same compact-subinterval sense.

Thus the limiting right-tail boundary is the parametrized curve
\[
 \boxed{(X,Y)=\left(\frac{s^2(2s-3)}{12},\frac{s(2-s)}4\right),\qquad3/2\le s\le2.}
\]
The left tail has the reflected profile, with its contributing angles of the form $w-\delta s+o(\delta)$.
\end{theorem}
\begin{proof}
The previous lemma confines every positive intercept to a fixed compact $s$ interval, where $Z_\delta\to V_*$ uniformly. The model function $V_*$ has unique global maximum $1/3$ at $s=2$. This proves both the floor-extent assertion and convergence of maximizing parameters. The supremum is attained at an interior angle for small $\delta$ because it is positive and its possible parameters are separated from the two original interval endpoints.

The scaled right-tail height is exactly
\[
 H_\delta(X)=\left[\sup_{0<s<w/\delta}(Z_\delta(s)-X)\,\delta\cot(\delta s)\right]_+.
\]
For $X$ in a compact subinterval of $(0,1/3)$, a positive contribution can come only from $s\le S$. It also cannot come from arbitrarily small $s$: the uniform estimate $Z_\delta(s)\le C\sqrt s+Cs^2$ is smaller than the positive lower bound for $X$ when $s$ is small enough. Thus the relevant parameters lie in a common compact subset of $(0,\infty)$.

On that set the expression converges uniformly to $(V_*(s)-X)/s$. Its positive maximum is attained at the unique parameter $s(X)$ computed in \texttt{transition-model.tex}, and has height $V_*'(s(X))=s(X)/2-s(X)^2/4$. Compactness and uniqueness of the maximizer give uniform convergence of all maximizing parameters when $X$ stays in a compact subinterval. Reflection supplies the final statement.
\end{proof}

\begin{theorem}[The two endpoint areas and the uncounted core]
For arbitrary maximizing caps,
\[
 \frac{|R_\delta|}{\delta^5}\longrightarrow\frac{229}{7680},\qquad
 \frac{|L_\delta|}{\delta^5}\longrightarrow\frac{229}{7680},
\]
and
\[
 |N_w(K_\delta)|-I(\gamma_\delta)-|R_\delta|-|L_\delta|=o(\delta^5).
\]
\end{theorem}
\begin{proof}
The model lower-bound proof gives the lower limit $229/7680$ for each tail. Localization of the quadratic gap gives lower limit $229/3840$ for that gap, because each of the two limiting derivative profiles is $v_*$. The sum of these three lower limits is $229/1920$, already the exact total deficit coefficient. The uncounted core correction is nonnegative by the disjoint core/tail decomposition. No one of these four nonnegative contributions can retain an additional positive limiting amount without contradicting the proved total value limit. This proves all three assertions. The argument controls area even though no uniform height statement at $X=0$ has been assumed.
\end{proof}

\begin{remark}[What cannot be inferred from this result]
The explicit curve is a rescaled limiting boundary, not the exact boundary at a fixed larger angle. The $O(\delta^6)$ core estimate was proved for the recovery sequence, not for every maximizing cap; the assertion here for maximizing caps is only $o(\delta^5)$. No uniform height convergence at $X=0$, no endpoint support displacement coefficient, and no exact larger-angle uniqueness theorem is claimed.
\end{remark}

\paragraph{Consequence for continuation.}
The exposed right-tail walls approach angle zero as $\omega\downarrow1$, while the reflected left-tail walls approach the final angle. A proposed fixed-angle contact description must reproduce these scales and the two limiting parameters $3/2$ and $2$. A contact diagram guessed solely from the right-angle endpoint does not supply that verification. No CI or compilation was run.
\end{document}
